\textbf{Occam's Razor} is the principle that, among competing explanations, the \textbf{simplest one} should be preferred. \vspace{7pt}

In machine learning, this refers to:
\begin{enumerate}
    \item Take always \textbf{simpler hypothesis} that fit the data.
    \item Avoid unnecessarily complex classifiers.
\end{enumerate}
Specifically, it states that any hypothesis that is consistent with a sufficiently larger set of training data is very likely to generalize well to unseen data, 
provided it comes from a relatively small hypothesis space. The main intuition can be divided into:
\begin{enumerate}
    \item A hypothesis that is \textbf{too complex} may \textbf{overfit} the data, such classifier provides good results only on training data.
    \item Simpler hypotheses are more likely to \textbf{generalize}, performing well also on unseen data.
\end{enumerate}